% =====================================================================
%  Chapter 10 — Principle: Rigor (Refuse the Almost Right)
%  Map: research/chapter-map.md (Chapter 10)
%  Shape: idea → origin → cost → what it opens up. Ideas, not how-to.
%  Rule: the author is never quoted in body text; short quotes only in
%  footnotes. Full statements live in interviews/.
% =====================================================================
\chaptersubtitle{Refuse the Almost Right}
\chapter{Rigor}
\label{ch:refuses}

A conventional compiler is generous. It accepts nearly anything that is
grammatical, grumbles about some of it in warnings, and then quietly repairs
the rest: it removes work that does nothing, merges work done twice, and
tidies arithmetic. Silica's compiler starts from the opposite assumption.
Code that is almost right is not right, and the compiler should say so rather
than cover for it.

\section{Refuse, don't repair}

Silica treats the patterns an optimizer usually cleans up---values computed and
never used, the same work done twice, arithmetic that could be simpler---as
mistakes in the program rather than chances for the compiler to look clever.
They are errors, not warnings, and they cannot be switched off. A warning is
something people learn to ignore. A repair the programmer never sees makes the
program's real behaviour harder to predict. Silica would rather the program say
what it means. The stance has a history: Barney's distrust of optimizers began
with seeing how much of their work goes into silently translating an old machine
model onto new chips (Chapter~\ref{ch:hidden}).

The beginner book frames refusal as a kindness: a program that is stopped cannot
quietly do the wrong thing.\footnote{\src{silica/docs/learn-programming.md},
line 206: ``That is a gift.''}

% TODO-VERIFY: Only one of these refusals (unused values) is implemented; the
%   rest are specified intent. Keep the text accurate as they land.
% TODO-INTERVIEW: Is the gap between the published rules and the implemented
%   one deliberate, a statement of intent?

\section{Let the machine do the risky work}

Rigor also means knowing which jobs people should not do by hand. Some routine
work is tedious, easy to get subtly wrong and dangerous when wrong: checking
data that arrives from a network, copying a message safely from one actor to
another, deciding when memory may be released. In most languages programmers
write that machinery themselves, and it is where security holes are commonly
born. In Silica the programmer does not write it at all. The program states
what kind of data it expects, and the compiler produces the checking, copying
and releasing from that statement when it builds the program.

\section{The whole program, seen whole}

Silica treats an application as one program, examined as a whole, not as a pile
of separately built pieces stitched together at the end. Libraries arrive as
source, not as opaque binaries. Seeing everything at once lets the compiler make
promises about the whole application that no piecemeal process could. Where the
current compiler cuts corners to save memory, the project treats that as a
temporary limitation to be fixed in the compiler, never as a reason to change
the language.

% TODO-INTERVIEW: Does the whole-program idea rule out ever updating a running
%   system in place, which Erlang is famous for?

\section{Programs as proofs}

Silica's type system is designed with a deep idea from logic in mind: that a
type can be read as a claim and a program as its proof. The project is candid
that this is a direction rather than a destination; no machine-checked proofs
come out of the compiler yet. Where the direction leads is a reversal of the
usual order of work: a proof of correctness written first, by a human who answers
for it, and code written afterwards, possibly by a model, to satisfy it
(Chapter~\ref{ch:memory}).

\section{Evidence over assertion}

Every claim the compiler makes about itself is tested. Tens of thousands of
small programs, each with its expected result, decide whether a change is
finished, and it is not finished until they pass on every kind of machine it
might affect. The compiler itself is the largest test. Written in Silica, it
compiles itself, and it is accepted only when the compiler it produces, compiling
itself again, gives back exactly the same result, byte for byte. That was first
achieved on 12~September 2026.

% TODO-INTERVIEW: Take me to the day of the first byte-for-byte match.

\section{What it costs}

A strict compiler is a demanding companion. Code that would run elsewhere is
refused here until it is clean. Familiar escape hatches are missing. And the
standard is costly to uphold: tens of thousands of tests, run on every platform,
for every change.

\section{What it opens up}

\begin{itemize}
  \item \textbf{Performance you can trust.} When nothing is rewritten behind
        the programmer's back, what is written is what runs.
  \item \textbf{Fewer places for attacks to start.} The routines where
        hand-written mistakes most often become security holes---checking
        incoming data, copying messages, releasing memory---are generated by
        the compiler from what the program declares, the same way every time,
        instead of being written anew by each programmer.
  \item \textbf{Builds you can verify.} A compiler that reproduces itself
        exactly is a foundation for checking that the software you run is the
        software you meant---a growing concern for anyone worried about
        tampered tools.
  \item \textbf{A road to proof.} A language shaped for proof from the start
        can grow into one that proves things, without a redesign.
\end{itemize}

% TODO-INTERVIEW: What does a compiler that refuses make possible for teams,
%   or for students, that a generous one does not?
